\section{Proposed Methodology}
\label{sec:methodology}

\subsection{Data Pipeline, Preprocessing, and Leakage Controls}
The empirical pipeline processes C-MAPSS subsets FD001, FD002, and FD004 (totaling 609 engines). All data partitioning is executed strictly at the independent engine cluster level: approximately 50\% training, 20\% calibration, 10\% validation, and 20\% held-out test engines (manifests preserved in \texttt{data/splits/}).

Preprocessing uses training-only z-score normalization. Means and population standard deviations are estimated with \texttt{ddof=0}. An exactly zero standard deviation is replaced by an effective scale of 1.0, while the mean is still subtracted:
\begin{equation}
x_{\text{normalized}} = \frac{x - \mu_{\text{train}}}{s_{\text{train}}}
\end{equation}
where $s_{\text{train}} = \sigma_{\text{train}}$ when $\sigma_{\text{train}} > 0$, and $s_{\text{train}} = 1.0$ otherwise.

Dataset-specific processing is structured as follows:
\begin{itemize}
  \item \textbf{FD001}: Single nominal operating condition; global training-fitted z-score normalization across all qualified features.
  \item \textbf{FD002 and FD004}: Operational flight settings are standardized using training data; six operating regimes ($C_1$--$C_6$) are learned from standardized training settings using $K$-Means clustering ($k=6$, $\text{random\_state}=2026$, $\text{n\_init}=10$). Canonical regime labels are assigned based on lexicographic centroid ordering. Sensor normalization parameters are fitted separately within each training regime.
\end{itemize}
Pre\begin{figure}[htbp]
\centering
\resizebox{\columnwidth}{!}{% Architecture Diagram for IntelliTwin Pipeline
% Strict partition isolation, split counts, and offline evaluation boundaries
\begin{tikzpicture}[
    node distance=1.2cm and 0.8cm,
    box/.style={rectangle, draw=black!70, rounded corners=2pt, minimum width=2.4cm, minimum height=0.9cm, align=center, font=\footnotesize},
    databox/.style={rectangle, draw=blue!70, fill=blue!5, rounded corners=2pt, minimum width=2.6cm, minimum height=1.0cm, align=center, font=\footnotesize},
    procbox/.style={rectangle, draw=teal!70, fill=teal!5, rounded corners=2pt, minimum width=2.6cm, minimum height=1.0cm, align=center, font=\footnotesize},
    evalbox/.style={rectangle, draw=purple!70, fill=purple!5, rounded corners=2pt, minimum width=2.6cm, minimum height=1.0cm, align=center, font=\footnotesize},
    arrow/.style={-{Stealth[length=2.5mm]}, thick, draw=black!70},
    dashedarrow/.style={-{Stealth[length=2.5mm]}, thick, dashed, draw=red!70}
]

% 1. Dataset Partition Layer
\node[databox] (split) {
    \textbf{Engine-Level Disjoint Split}\\
    \scriptsize FD001: 50 / 20 / 10 / 20\\
    \scriptsize FD002: 130 / 52 / 26 / 52\\
    \scriptsize FD004: 124 / 50 / 25 / 50\\
    \tiny (Train / Cal / Val / Test)
};

% 2. Partition Nodes
\node[procbox, right=1.2cm of split, yshift=2.0cm] (train) {
    \textbf{Training Partition}\\
    \scriptsize Training-only z-score fit\\
    \scriptsize 6 Regimes ($k$-means)\\
    \scriptsize GBR / LSTM fitting\\
    \scriptsize Target: $\min(\text{rul}, 125)$
};

\node[procbox, right=1.2cm of split, yshift=0.7cm] (val) {
    \textbf{Validation Partition}\\
    \scriptsize Validation-only selection\\
    \scriptsize Grid search ($\gamma, \beta$)\\
    \scriptsize Control parameter selection\\
    \scriptsize No test leakage
};

\node[procbox, right=1.2cm of split, yshift=-0.7cm] (cal) {
    \textbf{Reserved Calibration}\\
    \scriptsize Engine-weighted quantiles\\
    \scriptsize Finite mass: $\frac{m}{m+1}$\\
    \scriptsize Mass at $+\infty$: $\frac{1}{m+1}$\\
    \scriptsize Regime support $m_g \ge 10$
};

\node[evalbox, right=1.2cm of split, yshift=-2.2cm] (test) {
    \textbf{Held-Out Test Partition}\\
    \scriptsize Blended held-out trajectories\\
    \scriptsize \textbf{NO FEEDBACK} to selection\\
    \scriptsize Evaluated offline only
};

% 3. Decision & Simulator Layer
\node[procbox, right=1.2cm of val, yshift=-0.7cm] (policy) {
    \textbf{Frozen Decision Policy}\\
    \scriptsize Input: $L(t) = \max(0, p(t) - q a(p))$\\
    \scriptsize Rule: Maintain iff $L(t) \le 1.0$\\
    \scriptsize Evaluation: cycles $30 \dots T-1$\\
    \scriptsize Zero lead-time intervention
};

\node[evalbox, right=1.2cm of policy] (sim) {
    \textbf{Offline Simulation Replay}\\
    \scriptsize Synthetic SCU pricing:\\
    \scriptsize $\CPM=1.0, \CFAIL=5.0, \CWASTE=0.20$\\
    \scriptsize True RUL \& failure time\\
    \scriptsize available \textbf{only} offline
};

% Connections
\draw[arrow] (split.east) -- ++(0.5,0) |- (train.west);
\draw[arrow] (split.east) -- ++(0.5,0) |- (val.west);
\draw[arrow] (split.east) -- ++(0.5,0) |- (cal.west);
\draw[arrow] (split.east) -- ++(0.5,0) |- (test.west);

\draw[arrow] (train.east) -| ++(0.4,-1.2) -| (policy.north);
\draw[arrow] (val.east) -| ++(0.4,-0.3) -- (policy.west);
\draw[arrow] (cal.east) -| ++(0.4,0.3) -- (policy.west);
\draw[arrow] (test.east) -| ++(0.5,0.7) -| (sim.south);
\draw[arrow] (policy.east) -- (sim.west);

% Anti-leakage boundary
\draw[dashed, red!80, line width=0.8pt] ($(test.north west)+(-0.2,0.2)$) rectangle ($(sim.south east)+(0.2,-0.2)$);
\node[red!80, font=\tiny\bfseries, above right] at ($(test.north west)+(-0.2,0.2)$) {Strict Offline Evaluation Boundary (No Feedback)};

\end{tikzpicture}
}
\caption{IntelliTwin pipeline architecture, split counts, and strict offline evaluation boundary.}
\label{fig:architecture}
\end{figure}

Pre-processing executes strictly prior to window construction. The pipeline enforces a frozen 17-feature selection supported by the recorded qualification procedure:
\begin{center}
\texttt{setting1, setting2, setting3, s2, s3, s4, s7, s8, s9, s11, s12, s13, s14, s15, s17, s20, s21}
\end{center}
Seven sensors ($s_1, s_5, s_6, s_{10}, s_{16}, s_{18}, s_{19}$) are excluded under this recorded qualification procedure; we do not claim they are mathematically constant or universally uninformative across all datasets and regimes.

\begin{figure}[htbp]
\centering
\resizebox{\columnwidth}{!}{% Causal Window and Target Construction Diagram
% Illustrates 30-cycle lookback, 1-29 warm-up, dual target attachment, and 30x17 vs 510 representation
\begin{tikzpicture}[
    scale=0.95,
    cyclebox/.style={rectangle, draw=black!60, minimum width=0.5cm, minimum height=0.6cm, align=center, font=\tiny},
    warmupbox/.style={rectangle, draw=gray!50, fill=gray!15, minimum width=0.5cm, minimum height=0.6cm, align=center, font=\tiny},
    windowbox/.style={rectangle, draw=blue!80, fill=blue!10, minimum width=0.5cm, minimum height=0.6cm, align=center, font=\tiny},
    endpointbox/.style={rectangle, draw=red!80, fill=red!15, minimum width=0.5cm, minimum height=0.6cm, align=center, font=\tiny\bfseries},
    targetbox/.style={rectangle, draw=purple!80, fill=purple!10, rounded corners=2pt, minimum width=2.4cm, minimum height=0.7cm, align=center, font=\scriptsize},
    repbox/.style={rectangle, draw=teal!80, fill=teal!10, rounded corners=2pt, minimum width=2.5cm, minimum height=0.8cm, align=center, font=\scriptsize},
    arrow/.style={-{Stealth[length=2mm]}, thick, draw=black!70}
]

% Timeline coordinates
% Warm-up cycles 1, 2, ..., 28, 29
\node[warmupbox] (c1) at (0,0) {1};
\node[warmupbox] (c2) at (0.6,0) {2};
\node[font=\tiny] at (1.3,0) {$\cdots$};
\node[warmupbox] (c29) at (2.0,0) {29};

% Window cycles t-29 ... t
\node[windowbox] (ct29) at (3.0,0) {\tiny $t{-}29$};
\node[windowbox] (ct28) at (3.8,0) {\tiny $t{-}28$};
\node[font=\tiny] at (4.6,0) {$\cdots$};
\node[windowbox] (ct1) at (5.4,0) {\tiny $t{-}1$};
\node[endpointbox] (ct) at (6.2,0) {$t$};

% Next cycles
\node[cyclebox] (ctplus) at (7.0,0) {\tiny $t{+}1$};
\node[font=\tiny] at (7.7,0) {$\cdots$};
\node[cyclebox] (cT) at (8.4,0) {$T$};

% Labels and Braces
\draw[thick, decoration={brace, mirror, raise=4pt}, decorate, draw=gray!70] 
    ($(c1.south west)+(0,-0.05)$) -- ($(c29.south east)+(0,-0.05)$) 
    node[pos=0.5, below=8pt, font=\scriptsize, align=center] {Warm-up Horizon (Cycles 1--29)\\ \tiny No prediction generated; causal buffer filling};

\draw[thick, decoration={brace, raise=4pt}, decorate, draw=blue!80] 
    ($(ct29.north west)+(0,0.05)$) -- ($(ct.north east)+(0,0.05)$) 
    node[pos=0.5, above=8pt, font=\scriptsize, align=center] {\textbf{Causal Observation Window} (Length = 30 cycles, Stride = 1)\\ \tiny Strictly within engine boundaries; zero future cycle leakage};

% Feature Representation Nodes below
\node[repbox] (lstm_rep) at (3.5, -2.0) {
    \textbf{LSTM Representation}\\
    \scriptsize Matrix: $30 \times 17$\\
    \tiny 30 time steps $\times$ 17 features
};

\node[repbox] (gbr_rep) at (6.8, -2.0) {
    \textbf{GBR Representation}\\
    \scriptsize Vector: $510$ features\\
    \tiny C-order row-major flattening
};

\draw[arrow] (4.6, -0.4) -- (lstm_rep.north);
\draw[arrow] (5.4, -0.4) -- (gbr_rep.north);

% Target nodes attached to endpoint t
\node[targetbox] (train_target) at (9.6, 1.0) {
    \textbf{Training Target}\\
    \scriptsize $y_{\text{train}} = \min(\text{rul}_{\text{true}}, 125)$ cycles\\
    \tiny Linear piecewise modeling convention
};

\node[targetbox] (eval_target) at (9.6, -0.6) {
    \textbf{Evaluation Target}\\
    \scriptsize $y_{\text{eval}} = \text{rul}_{\text{true}} = T - t$ cycles\\
    \tiny Uncapped true remaining cycles
};

\draw[arrow, draw=purple!80] (ct.north) |- (train_target.west);
\draw[arrow, draw=purple!80] (ct.south) |- (eval_target.west);

\end{tikzpicture}
}
\caption{Causal observation window ($w=30$), feature representations (LSTM $30 \times 17$ vs.\ GBR 510), and target definitions.}
\label{fig:causal_window}
\end{figure}

Causal temporal windowing extracts chronological sequences of length $W=30$ cycles with stride 1, strictly no padding, and no crossing of engine boundaries. Predictions start at cycle 30. In training splits, target RUL is capped at $\min(\text{rul\_true}, 125)$ cycles as a linear piecewise modeling convention. Point predictions are evaluated as $p = \max(0, \text{raw\_prediction})$ clipped below at zero without upper clipping at 125.

\subsection{Frozen Base Predictor Matrix}
Two base predictor families receive identical 30-cycle causal information horizons:
\begin{itemize}
  \item \textbf{Histogram Gradient Boosting Regressor (GBR)}: \texttt{sklearn.ensemble.HistGradientBoostingRegressor} operating on C-order row-major flattened $30 \times 17$ windows (510 features). Winning configurations selected on validation data across seeds 2026--2027: \texttt{gbr\_lr010\_leaf15\_l2\_0} for FD001, and \texttt{gbr\_lr010\_leaf31\_l2\_0} for FD002 and FD004.
  \item \textbf{Unidirectional LSTM}: PyTorch neural network comprising one unidirectional recurrent layer (\texttt{input\_size=17}, \texttt{hidden\_size=64}, \texttt{num\_layers=1}, \texttt{batch\_first=True}, \texttt{bidirectional=False}, internal recurrent \texttt{dropout=0.0}), extracting final sequence output followed by \texttt{nn.Dropout(p=0.20)} and \texttt{nn.Linear(64, 1)}. Optimization uses AdamW with learning rate $\eta = 0.001$, weight decay $10^{-4}$, batch size 256, betas $(0.9, 0.999)$, and epsilon $10^{-8}$. Early stopping is governed by validation uncapped-RUL RMSE. Winning configuration: \texttt{lstm\_h64\_lr0010}.
\end{itemize}
Each selected configuration is trained across five final research seeds (\texttt{2026, 2027, 2028, 2029, 2030}), yielding 30 frozen checkpoints cataloged in \texttt{experiments/base\_models/model\_registry.json}. We do not characterize the five GBR seeds as independent experiments, nor assert universal deterministic behavior beyond the recorded runs.

\subsection{Hierarchical Engine-Weighted Conformal Calibration}
Standard pooled conformal prediction introduces survivor bias in industrial run-to-failure telemetry because long-lived engines contribute more sliding windows than short-lived units. IntelliTwin implements hierarchical engine-weighted quantiles motivated by Lee, Barber, and Willett~\cite{lee2023hierarchical}.

IntelliTwin uses an engine-weighted conformal construction motivated by hierarchical conformal prediction. Formal coverage results require the hierarchical sampling and exchangeability assumptions stated in the underlying theory. Equal engine weighting does not establish these assumptions for ordered, lifetime-stratified degradation trajectories. This study therefore evaluates marginal and subgroup coverage empirically. It does not claim guaranteed coverage for each engine, lifecycle stage, operating regime, complete trajectory, or adaptively selected maintenance time. Nominal 90\% coverage is a target, not a guaranteed realized coverage proportion. Maintenance outcomes are evaluated separately under the frozen simulator.

For FINAL calibration with $m$ contributing engines, where engine $i$ contains $n_i$ valid endpoints, the empirical distribution function is:
\begin{equation}
F_{\text{final}}(q) = \frac{1}{m + 1} \sum_{i=1}^m \frac{1}{n_i} \sum_{t=1}^{n_i} \mathbb{I}(\text{score}_{it} \le q)
\end{equation}
The mass of each finite endpoint is $\frac{1}{(m+1)n_i}$. The total finite mass is $\frac{m}{m+1}$, with the remaining probability mass $\frac{1}{m+1}$ placed at $+\infty$. The conformal threshold $q^*$ for nominal coverage $1 - \alpha$ is the smallest observed score whose CDF reaches $1 - \alpha$. If no finite score reaches it, $q^* = +\infty$. (In contrast, normalized development empirical CDFs use $1/m$ weighting when selecting configurations via leave-one-engine-out cross-validation).

\subsection{Formulations of Evaluated Conformal Methods}
Let $p \ge 0$ denote non-negative predicted RUL in cycles, $y = \trueRUL \ge 0$ denote uncapped ground-truth RUL, and $g \in \{C_1, \dots, C_6\}$ denote the observed operational regime.

\subsubsection{Global Hierarchical Conformal Prediction (\texttt{global\_hcp})}
Uses symmetric absolute residuals $\text{score}(p, y) = |y - p|$, fitting a single global threshold $q_{\text{global}}$, yielding $L = \max(0, p - q_{\text{global}})$ and $U = p + q_{\text{global}}$. The baseline interval is symmetric before lower clipping, with displayed width $p + q - \max(0, p - q)$, which equals $2q$ only when $p \ge q$.

\subsubsection{Condition-Aware Hierarchical Conformal Prediction (\texttt{regime\_hcp})}
Uses symmetric absolute residuals, fitting separate regime-specific thresholds $q_g$. Regime support requires $m_g \ge \max(10, \lceil 1/\alpha \rceil - 1)$ distinct contributing engines (10 engines for 80\% and 90\% coverage; 19 engines for 95\% coverage). If $m_g$ is insufficient, the method falls back deterministically to $q_{\text{global}}$ fitted for the same method and coverage level, recording fallback.

\subsubsection{Proposed Decision-Aware Asymmetric Conformal Prediction (\texttt{decision\_aware\_hcp})}
Modulates interval width and asymmetry via scaling functions:
\begin{align}
\sigma(p) &= \left(1 + \frac{p}{125}\right)^\gamma, \quad z(p) = \exp\left(-\frac{p}{20}\right) \\
a(p) &= \sigma(p)\exp(\beta z(p)), \quad b(p) = \sigma(p)\exp(-\beta z(p))
\end{align}
where $\gamma \in \{0, 1, 2\}$ and numerical $\beta \in \{-\ln 2, 0, +\ln 2\}$. Nonconformity scores and inverted intervals are:
\begin{equation}
\text{score}(p, y) = \max\left(\frac{p - y}{a(p)}, \; \frac{y - p}{b(p)}, \; 0\right)
\end{equation}
\begin{equation}
L = \max(0, \; p - q_g \cdot a(p)), \quad U = p + q_g \cdot b(p)
\end{equation}

\subsection{Mathematical Mechanics of Negative Beta ($\beta < 0$)}
At fixed $p$, $q$ and $\gamma$, negative $\beta$ decreases $a(p)$ and increases $b(p)$. The lower and upper endpoints are therefore nondecreasing relative to $\beta=0$. This does not establish earlier intervention or better maintenance utility because calibration thresholds are refitted and policy outcomes depend on the resulting trajectories.

Candidate \texttt{gamma2\_betam1} ($\gamma=2, \beta_{\text{code}}=-1, \beta=-\ln 2$) was selected as the optimal feasible configuration across all six mandatory cells and frozen under ADR-021.
